\documentclass[10pt,a4paper]{article}
\usepackage[margin=23mm]{geometry}
\usepackage{amsmath,amssymb,booktabs,array,hyperref}
\hypersetup{colorlinks=true,urlcolor=blue,linkcolor=blue}
\setlength{\parindent}{0pt}
\setlength{\parskip}{4pt}
\newcommand{\scl}{\operatorname{scl}}
\newcommand{\cl}{\operatorname{cl}}
\newcommand{\PSL}{\operatorname{PSL}}
\newcommand{\SL}{\operatorname{SL}}
\title{A rank-one lattice with unbounded stable commutator length\\
\large Counterexample to problem 00000001158}
\author{lizaixi01}
\date{24 October 2026}
\begin{document}
\maketitle
\vspace{-6mm}

\textbf{Scope and result.}
The printed conjecture asserts that a lattice has bounded stable commutator
length if and only if its rank is one, with a knot-complement exception.
Use the usual ambient real-rank convention for lattices in semisimple real
Lie groups. Put
\[
 \overline{\Gamma(2)}=\Gamma(2)/\{\pm I\}<\PSL_2(\mathbb R),\qquad
 \Gamma(2)=\ker\bigl(\SL_2(\mathbb Z)\longrightarrow\SL_2(\mathbb F_2)\bigr).
\]
This standard nonuniform lattice has ambient real rank one and is a free
group $F=\langle a,b\rangle$ of rank two. For $c=[a,b]$, the supplied Lean
proof establishes, for every integer $N\ge0$,
\begin{equation}\label{bounds}
             \frac{N}{12}\ \le\ \scl_F(c^N)\ \le\ N.
\end{equation}
Thus the actual finite values on the commutator subgroup are unbounded.
The example is not the group of a usual single knot in $S^3$: its
abelianization has a noncyclic quotient $\mathbb Z^2$, whereas such knot
groups have cyclic abelianization. The conjecture prints no cocompactness
restriction. No extended-infinite value outside $[F,F]$ is used.

\textbf{Published geometric identification.}
The exact geometric and topological identifications in this paragraph are
used from published standard models; they are \emph{not re-proved in Lean}.
Morris, Example 1.3.7, proves that $\SL_2(\mathbb Z)$ has finite Haar
covolume in $\SL_2(\mathbb R)$ by integrating over its modular fundamental
domain. Corollary 4.1.14(2) preserves the lattice property under finite
index. The kernel defining $\Gamma(2)$ has finite index, as Conrad explains
before Theorem 3.1. Quotienting both groups by the finite central subgroup
$\{\pm I\}$ preserves discreteness and the finite invariant quotient
measure, since the corresponding lattice quotient spaces agree.
Morris, Example 8.1.7(1) and Remark 8.1.9(2), give real rank
$\operatorname{rank}_{\mathbb R}\SL_2(\mathbb R)=1$ and invariance under
compact central quotients. Conrad, Remark 3.10 (pp.10--11), identifies
$\overline{\Gamma(2)}$ with a free group of rank two. Importantly,
$\Gamma(2)$ itself contains $-I$ and is not that free group.
Burde--Zieschang, Chapter 8A, p.103, gives the cyclic abelianization of
classical knot groups. The Lean theorem transports the algebraic result
through every group isomorphism; no geometric axiom or assigned lattice
or rank label is introduced.

\textbf{Genuine minimum and genuine stable limit.}
For an arbitrary group $G$ and a finite list $L=((u_i,v_i))$, define its
actual product $P(L)=\prod_i[u_i,v_i]$. The proof shows that all such
products form precisely Mathlib's native subgroup $[G,G]$; inverse
products are represented by reversing the list and swapping the pairs.
For every $x\in[G,G]$ the minimum
\[
  \cl_G(x)=\min\{|L|:P(L)=x\}
\]
therefore exists. It is implemented by \texttt{Nat.find} only after this
existence proof, with an actual minimizing witness and the inequality
against \emph{every} representing list. Concatenation proves
$\cl_G(xy)\le\cl_G(x)+\cl_G(y)$. The sequence
$u_n=\cl_G(x^n)$ is subadditive, and $u_n/n\ge0$. Mathlib's native
Fekete theorem then proves the genuine limit
\[
 \scl_G(x)=\lim_{n\to\infty}\frac{\cl_G(x^n)}n
 =\inf_{n\ge1}\frac{\cl_G(x^n)}n,
 \qquad 0\le\scl_G(x)\le\cl_G(x).
\]
The nonempty, bounded-below sequence supplies the infimum and its limit;
neither a desired norm nor a fallback infimum value is assigned.

\textbf{Actual Brooks count and its defect.}
For the native free group on two generators, let $w(g)$ be its actual
reduced word. Write $A=a^{-1}$ and $B=b^{-1}$. The integer
\[
 H(g)=\#\{ab\text{ occurrences in }w(g)\}
       -\#\{BA\text{ occurrences in }w(g)\}
\]
is constructed from adjacent-letter indicators. The proof identifies
the recursive sum with actual \texttt{List.zip}, \texttt{tail} and
\texttt{countP} occurrence counts. It proves $H(g^{-1})=-H(g)$.
For any two reduced words the proof constructs the full cancellation
decomposition $w(g)=UC$, $w(h)=C^{-1}V$, with $UV$ reduced. Each of
the three remaining word junctions changes the signed count by at most
one. Hence
\[
 |H(gh)-H(g)-H(h)|\le3,\qquad |H([g,h])|\le9.
\]
Induction on every commutator-product list gives
$|H(P(L))|\le12|L|$. These are proved bounds for the concrete count,
consistent with the standard counting construction in Calegari,
Definition 2.25 and Lemma 2.29; no quasimorphism is assumed.

\textbf{Unbounded finite stable lengths.}
The native word of $c$ is $abAB$. It is cyclically reduced, and the proof
identifies the exact reduced word of every $c^n$ with the concatenation
of $n$ copies of this word. Its actual occurrence count is $H(c^n)=n$.
Consequently, for every positive $n$ and every $N\ge0$,
\[
 Nn=H((c^N)^n)\le12\cl_F((c^N)^n),\qquad
 \frac{N}{12}\le\frac{\cl_F((c^N)^n)}n.
\]
Passing to the genuine stable limit proves the lower bound in
\eqref{bounds}. A list of $N$ copies of the pair $(a,b)$ represents
$c^N$, proving the upper bound. For every real $C$, choose a natural
$N>12C$; then $C<\scl_F(c^N)\le N$. Lean proves this quantified
witness statement and the failure of boundedness of the full range.
It also proves that genuine commutator lengths and stable limits are
preserved by every group isomorphism, by transporting all product lists
in both directions and using uniqueness of limits.

\textbf{The knot exception is excluded.}
The genuine exponent-sum homomorphism
$F\to(\mathbb Z^2,+)$ sends $a$ to $(1,0)$ and $b$ to $(0,1)$.
It is onto: $a^m b^n$ maps to $(m,n)$ for all integers $m,n$.
Mathlib's universal map through the native \texttt{Abelianization}
therefore gives a surjection $F^{\rm ab}\twoheadrightarrow\mathbb Z^2$.
The proof rules out a cyclic generator $(r,s)$ of $\mathbb Z^2$:
representing $(1,0)$ forces $s=0$, making $(0,1)$ impossible.
Thus $F^{\rm ab}$ is not cyclic. A final theorem excludes every group
isomorphism from $F$ to any group with cyclic native abelianization,
which applies to the published knot-group theorem above. This argument
does not address arbitrary link complements or knots in other ambient
manifolds.

\textbf{Formalization and reproducibility.}
The complete algebraic proof is in \texttt{lean/Main.lean}, namespace
\texttt{TLMC1158}, culminating in \texttt{complete\_counterexample}.
The geometric model identification, Haar covolume and Lie rank, and the
general knot-group theorem are the explicitly disclosed external
boundaries. The proof neither formalizes these identifications by an
extra axiom nor claims to prove a general classification of lattices or
knot groups. The target is only the printed rank-one boundedness
implication under the stated conventions.

The project pins Lean 4.33.0 and Mathlib commit
\texttt{db584cd6d46c92f209a44c0f1c829460d327499d}.
Run \texttt{lean --version}, then \texttt{lean/verify.sh} for the build,
ordinary exhaustive namespace axiom audit and official fresh-kernel
replay. Command logs, original job and resource observations, exact
source bindings and the file manifest accompany a verified submission.
Production timings do not establish a controlled efficiency gain.

\small
\begin{thebibliography}{4}
\bibitem{calegari} D. Calegari, \emph{scl}, MSJ Memoirs 20, Mathematical
Society of Japan, 2009, pp.13,18,21--22.
\url{https://www.math.uchicago.edu/~dannyc/books/scl/scl.pdf}.
\bibitem{conrad} K. Conrad, \emph{SL$_2(\mathbb Z)$}, Theorem 3.1 and
Remark 3.10, pp.7,10--11.
\url{https://kconrad.math.uconn.edu/blurbs/grouptheory/SL(2,Z).pdf}.
\bibitem{morris} D. W. Morris, \emph{Introduction to Arithmetic Groups},
arXiv:math/0106063v6, 2015, pp.15--16,46,180--181.
\url{https://arxiv.org/abs/math/0106063v6}.
\bibitem{knots} G. Burde and H. Zieschang, \emph{Knots}, second revised
and extended edition, de Gruyter, 2003, Chapter 8A, p.103.
\url{https://www.maths.ed.ac.uk/~v1ranick/papers/burdzies.pdf}.
\end{thebibliography}
\end{document}
